\subsection{Characterization of bounded measures}

PROPOSITION 12. --- For a measure \(\mu\) on a locally compact space \(X\) to be bounded (Ch. III, §1, No.~8), it is necessary and sufficient that \(X\) be an integrable set with respect to \(\mu\) (or, what comes to the same thing, that every finite constant function be integrable); in this case,

\[
\| \mu \| = | \mu | (\mathrm{X}) = \int d | \mu |.
\]

For, we have seen that \(|\mu|^{*}(X) = \|\mu\|\) (§1, No.~2); the proposition therefore follows from Prop. 10 of No.~6.

For every bounded measure \(\mu\) , we again say that \(\mu(\mathrm{X})\) is the total mass of \(\mu\) .

It follows from Th. 4 of No.~6 that if \(\mu\) is a bounded measure then, for every \(\varepsilon > 0\) , there exists a compact set \(K\) such that \(|\mu| (\mathbb{C}K) \leqslant \varepsilon\) .

PROPOSITION 13. --- Let \(\mu\) be a bounded measure on X. Let \(\mathfrak{B}\) be a filter base on \(\mathcal{L}_{\mathrm{F}}^{p}\) having the following properties:

\(1^{\circ}\) there exists a set \(\mathbf{M} \in \mathfrak{B}\) such that the functions \(\mathbf{f} \in \mathbf{M}\) are uniformly bounded on \(\mathbf{X}\) ;

\(2^{\circ}\) \(\mathfrak{B}\) converges uniformly on every compact subset of X to a function \(\mathbf{f}_0\) .

Under these conditions, \(\mathbf{f}_0\) belongs to \(\mathcal{L}_{\mathrm{F}}^{p}\) and \(\mathfrak{B}\) converges in mean of order \(p\) to \(\mathbf{f}_0\) .

We note first of all that if \(|\mathbf{f}(x)| \leqslant a\) for every \(x \in X\) and every function \(\mathbf{f} \in M\) , then also \(|\mathbf{f}_0(x)| \leqslant a\) for every \(x \in X\) . This being so, for every \(\varepsilon > 0\) there exist a compact set \(K\) such that \(|\mu|(\mathbb{C}K) \leqslant \varepsilon^p\) and a set \(N \in \mathfrak{B}\) such that, for every function \(\mathbf{f} \in N\) , \(|\mathbf{f}(x) - \mathbf{f}_0(x)| \leqslant \varepsilon(|\mu|(K))^{-1/p}\) for all \(x \in K\) . Now, we may write

\[
\mathbf {f} - \mathbf {f} _ {0} = (\mathbf {f} - \mathbf {f} _ {0}) \varphi_ {\mathrm{K}} + (\mathbf {f} - \mathbf {f} _ {0}) \varphi_ {\mathbf {C K}};
\]

it follows from the foregoing that if \(\mathbf{f} \in \mathrm{M} \cap \mathrm{N}\) then \(\mathrm{N}_p((\mathbf{f} - \mathbf{f}_0)\varphi_{\mathrm{K}}) \leqslant \varepsilon\) and \(\mathrm{N}_p((\mathbf{f} - \mathbf{f}_0)\varphi_{\mathbf{C}_{\mathrm{K}}}) \leqslant 2a\varepsilon\) , whence \(\mathrm{N}_p(\mathbf{f} - \mathbf{f}_0) \leqslant (2a + 1)\varepsilon\) , which proves the proposition.

COROLLARY. --- For a bounded measure \(\mu\) on \(X\) , every bounded continuous mapping \(\mathbf{f}\) of \(X\) into \(\mathbf{F}\) belongs to each of the \(\mathcal{L}_{\mathrm{F}}^{p}\) ( \(1 \leqslant p < +\infty\) ).

For every compact subset K of X, let \(M_{K}\) be the set of mappings of X into F of the form hf, where h is a continuous mapping of X into {[}0,1{]} equal to 1 on K and with compact support. It is clear that the sets \(M_{K}\) form a filter base B on \(L_{F}^{p}\) , that the functions belonging to \(M_{K}\) are uniformly bounded, and that B converges uniformly to f on every compact subset of X, whence the corollary.

In particular, the function f is integrable and its integral \(\int f d\mu\) is the limit with respect to B of the integrals \(\int h f d\mu\) .

We will obtain anew the Cor. of Prop. 13 as a consequence of a general criterion for integrability in §5, No.~6.

In the notations of Ch. III, §1, No.~2, \(|\mathbf{f}| \leqslant \| \mathbf{f} \| \cdot 1\) for every function \(\mathbf{f} \in \mathcal{C}^b(\mathrm{X};\mathrm{F})\) , whence, by the formulas (3) and (4) of No.~2,

\[
\mathrm{N} _ {p} (\mathbf {f}) \leqslant \| \mathbf {f} \| \cdot \mathrm{N} _ {p} (1) = \| \mathbf {f} \| \cdot \| \mu \| ^ {1 / p}.\tag{19}
\]

In particular, for \(p = 1\) , formula (5) of No.~2 yields

\[
\left| \int \mathbf {f} d \mu \right| \leqslant \| \mathbf {f} \| \cdot \| \mu \|,\tag{20}
\]

consequently the mapping \(f \mapsto \int f d\mu\) is continuous on the Banach space \(\mathcal{C}^{b}(X;F)\) ; its restriction to the closure \(\mathcal{C}^{0}(X;F)\) of \(\mathcal{K}(X;F)\) in \(\mathcal{C}^{b}(X;F)\) , that is, to the space of continuous functions tending to 0 at the point at infinity (Ch. III, §1, No.~2, Prop. 3), is therefore the extension by continuity of the integral to \(\mathcal{C}^{0}(X;F)\) .
% semantic-fragments: legacy-input-seam
\relax{}
